\begin{center}
    \section*{\fontsize{20}{20}\selectfont Chapter 3}
\end{center}
\vspace{10mm}
\section{Research Methodology}

This chapter presents the methodology developed to build and validate the graph-based neural network for brain tumor segmentation. The research follows a systematic pipeline comprising data preprocessing, graph construction, model architecture design, training, and rigorous validation. Each stage is designed to ensure reproducibility, eliminate data leakage, and enable fair comparison with baseline methods.

The proposed framework converts 3D multi-parametric MRI volumes into 2D graph representations, where superpixels become nodes and spatial adjacency defines edges. This representation reduces computational complexity while preserving essential tumor characteristics through carefully engineered node features. The GNN processes these graphs using message-passing layers that aggregate neighbourhood information, enabling both local texture analysis and global structural reasoning. All experimental procedures follow deterministic protocols with comprehensive integrity checks to ensure scientific rigour and clinical reliability.


\subsection{Proposed Framework}

The framework progresses from raw multi-parametric MRI input through preprocessing, graph construction with superpixel-based dimensionality reduction ($\geq$890$\times$ compression), feature engineering (15 domain-specific features), GraphSAGE training, patient-level cross-validation, ensemble aggregation, efficiency benchmarking, and comprehensive validation. Each phase incorporates integrity checks to prevent data leakage.

The framework consists of eight interconnected phases that transform raw MRI data into validated tumor segmentation predictions.

\textbf{Phase 1: Data Preprocessing} standardizes the BraTS 2021 dataset through skull-stripping, resampling to isotropic 1mm$^3$ voxel spacing, Z-score normalization per modality, and patient-level stratified 5-fold cross-validation split generation. This prevents data leakage.

\textbf{Phase 2: Graph Construction} converts preprocessed 3D volumes into 2D graph representations. For each patient, we extract 200 axial slices using adaptive selection to focus on tumor-containing regions. The SLIC algorithm segments each slice into superpixels (compactness parameter = 10). Each superpixel becomes a graph node with 15 engineered features that capture intensity statistics, spatial information, and shape descriptors. Graph edges connect spatially adjacent superpixels based on boundary sharing.

\textbf{Phase 3: Feature Engineering} extracts domain-specific features without ground-truth leakage. The 15 features include intensity statistics from all four MRI modalities (T1, T1CE, T2, FLAIR), spatial coordinates, and geometric properties. Critically, we don't use any tumor label information during feature extraction—that would be cheating and would lead to data leakage.

\textbf{Phase 4: Model Training} employs a 5-layer GraphSAGE architecture with 256 hidden dimensions. We train using Adam optimizer (learning rate 0.001), binary cross-entropy loss with class weights to handle imbalance, batch size 32, and 100 epochs with early stopping. Training uses deterministic seed (42) and CUDA deterministic mode for reproducibility.

\textbf{Phase 5: Cross-Validation} implements patient-level stratified 5-fold splits, ensuring no patient appears in both training and validation sets. Each fold is trained independently with identical hyperparameters, enabling robust performance estimation and statistical validity assessment.

\textbf{Phase 6: Ensemble Methods} combine predictions from all 5 folds using soft voting (averaging sigmoid outputs) to leverage model diversity and improve robustness. The ensemble consistently yields a 2--3\% absolute improvement in Dice over individual fold models.

\textbf{Phase 7: Efficiency Benchmarking} compares inference speed, memory consumption, and model size against 3D U-Net baseline under identical hardware conditions. We measure per-patient latency, GPU memory usage, and parameter count.

\textbf{Phase 8: Integrity Validation} conducts 15 comprehensive audits to ensure research integrity and reproducibility: (1) feature dimension verification (15 features per node), (2) patient leakage detection across splits, (3) label distribution analysis across folds, (4) normalization consistency checks (Z-score per modality), (5) seed reproducibility verification (deterministic outputs), (6) graph node count validation ($\sim$3,900 nodes per patient on average, range 2,300--6,000), (7) edge connectivity verification (planar graph structure), (8) cross-fold data isolation audit (no shared patients), (9) ground truth label integrity (binary mask validity), (10) MRI modality completeness check (all 4 modalities present), (11) slice count validation (155 axial slices per BraTS patient, with $\sim$138 active slices after empty-slice filtering), (12) superpixel boundary consistency, (13) training/validation/test split ratio verification (900/100/251 patients, 72\%/8\%/20\%), (14) model checkpoint consistency, and (15) prediction output format verification.


\subsection{Data Set Analysis}

This research utilizes the Brain Tumor Segmentation (BraTS) 2021 dataset, the most comprehensive and widely adopted benchmark for evaluating automated brain tumor segmentation algorithms.

\textbf{Dataset Composition:} The BraTS 2021 training set contains 1,251 multi-institutional glioma cases collected from 19 institutions worldwide. Each case includes four co-registered MRI modalities: native T1-weighted (T1), T1-weighted with contrast enhancement (T1CE), T2-weighted (T2), and Fluid Attenuated Inversion Recovery (FLAIR). All images have been pre-processed by BraTS organizers to normalize orientations, resample to 1mm$^3$ isotropic resolution, and skull-strip to focus on brain tissue.

\textbf{Ground Truth Annotations:} Expert neuroradiologists manually segmented three nested tumor sub-regions: (1) enhancing tumor (ET), (2) tumor core (TC = ET + necrotic core), and (3) whole tumor (WT = TC + peritumoral edema). For this research, we focus on binary segmentation (tumor vs. non-tumor) by combining all tumor regions into a single positive class.

\textbf{Data Characteristics:} Glioma tumors exhibit high heterogeneity in appearance, size (ranging from small focal lesions to large infiltrative masses), location (various brain regions), and intensity patterns across modalities. Tumors typically occupy only 5-10\% of total brain volume, creating significant class imbalance. The multi-institutional nature introduces scanner variability (different field strengths, acquisition protocols, vendors) that tests generalization capability.

\textbf{Patient-Level Splitting:} Critical for valid medical AI research, we implement patient-level stratified 5-fold cross-validation where each patient's complete 3D volume appears in exactly one fold. This prevents slice-level leakage where slices from the same patient could appear in both training and testing, artificially inflating performance. Stratification ensures balanced tumor distribution across folds.

\begin{table}[htbp]
\caption{BraTS 2021 Dataset Statistics and Cross-Validation Split}
\begin{center}
\begin{tabularx}{0.8\textwidth} { 
  | >{\raggedright\arraybackslash}X 
  | >{\centering\arraybackslash}X 
  | }
\hline
\textbf{Characteristic} & \textbf{Value}\\
\hline
Total Patients & 1,251\\
\hline
MRI Modalities & 4 (T1, T1CE, T2, FLAIR)\\
\hline
Image Resolution & 240 $\times$ 240 $\times$ 155\\
\hline
Voxel Spacing & 1 mm$^3$ (isotropic)\\
\hline
Training Patients (per fold) & 900 (72\%)\\
\hline
Validation Patients (per fold) & 100 (8\%)\\
\hline
Test Patients (per fold) & 251 (20\%)\\
\hline
Average Tumor Volume & 5-10\% of brain\\
\hline
Ground Truth Labels & Binary (tumor/non-tumor)\\
\hline
Cross-Validation Folds & 5 (patient-level stratified)\\
\hline
Total Graph Nodes (per patient) & $\sim$10,000 superpixels\\
\hline
Node Features & 15 (intensity + spatial + shape)\\
\hline
\end{tabularx}
\label{tab1}
\end{center}
\end{table}

\subsubsection{Data Collection \& Preprocessing}

\textbf{Data Acquisition:} The BraTS 2021 dataset was obtained through official registration with the BraTS challenge organizers. All data is de-identified and approved for research use, complying with institutional review board (IRB) requirements and patient privacy regulations.

\textbf{Preprocessing Pipeline:} Although BraTS provides pre-processed images, we apply additional standardization:

\begin{itemize}
    \item \textbf{Skull-Stripping:} Remove non-brain tissue to focus computation on relevant regions (already performed by BraTS, verified in our pipeline)
    
    \item \textbf{Intensity Normalization:} Apply Z-score normalization independently for each MRI modality to standardize intensity distributions across patients: $x_{norm} = (x - \mu) / \sigma$ where $\mu$ and $\sigma$ are computed per modality per patient
    
    \item \textbf{Slice Selection:} Extract 200 axial slices per patient using adaptive sampling that prioritizes tumor-containing regions while maintaining anatomical coverage
    
    \item \textbf{Quality Control:} Validate data integrity through automated checks for missing files, corrupted images, inconsistent dimensions, and ground truth label validity
\end{itemize}

\subsubsection{Feature Engineering and Selection}

The success of graph-based segmentation depends on node features that capture tumor-discriminative information without introducing ground-truth leakage. We engineered 15 features per superpixel node, organized into three categories:

\textbf{1. Multi-Modal Intensity Statistics (8 features):} For each of the four MRI modalities (T1, T1CE, T2, FLAIR), two statistical descriptors are computed over all voxels within each superpixel:
\begin{itemize}
    \item Mean intensity: captures the average signal strength per modality
    \item Standard deviation: measures local intensity variation per modality
\end{itemize}

Glioma tumors exhibit characteristic intensity patterns: hypo-intense in T1, hyper-intense in T2/FLAIR, and rim enhancement in T1CE. Mean and standard deviation per modality enable the GNN to learn these discriminative multi-modal signatures, yielding 8 intensity features in total ($2 \times 4$ modalities).

\textbf{2. Spatial and Area Features (4 features):} Superpixel area (absolute pixel count), normalized area (fraction of total slice area), and normalized centroid coordinates ($y$, $x$) provide spatial context, enabling the model to learn anatomical priors about typical tumor locations and region scale.

\textbf{3. Shape and Texture Features (3 features):} Perimeter (boundary pixel count via morphological erosion), compactness (circularity measure, $4\pi A / P^2$, where $A$ is area and $P$ is perimeter), and intensity range (peak-to-peak signal across all four modalities within the superpixel) capture morphological and textural properties that distinguish irregular tumor boundaries from smooth healthy tissue.

\textbf{Critical Design Decision:} No ground-truth label information is included in the feature set. Initial experiments included a ``tumor\_ratio'' feature (the percentage of tumor-labelled pixels within each superpixel), which achieved 98\% accuracy. However, this constitutes severe data leakage, as ground-truth labels are used as model input. Following removal of this feature and retraining, the validated performance of 90.39\% was obtained, ensuring fair evaluation and genuine generalisation capability.


\subsection{Algorithm/ Model Analysis}

\subsubsection{Graph Construction Algorithm}

The graph construction process converts each 3D MRI volume into a 2D graph representation suitable for GNN processing. The procedure consists of five systematic steps:

\textbf{Step 1: Slice Extraction} - Select 200 axial slices from the 3D volume (indices 20-220 of the 240-slice stack) to focus on brain tissue while maintaining computational efficiency.

\textbf{Step 2: Superpixel Generation} - Apply SLIC (Simple Linear Iterative Clustering) algorithm to each slice with parameters: n\_segments adaptively set based on brain mask size (up to 200 per slice, computed as brain\_pixels / 200), compactness=0.1 (favouring irregular, image-content-following superpixels over spatially uniform blocks), and Gaussian smoothing $\sigma=0.3$ applied prior to clustering. SLIC groups similar pixels into perceptually coherent regions that often align with anatomical structures.

\textbf{Node Count Design Choice:} The adaptive SLIC configuration yields approximately \textbf{46 superpixels per active slice} on average (range 10--200 depending on brain mask coverage), with $\sim$138 active slices per patient, resulting in roughly \textbf{3,900 total nodes per patient} (measured: mean 3,909 $\pm$ 821, range 2,347--6,041 across 50 patients). This design balances computational efficiency with spatial resolution — denser superpixel sampling (e.g., 200 nodes/slice on all 155 slices yielding $\sim$31K nodes) would preserve finer texture detail but increase graph processing cost by $\sim$8$\times$, while sparser sampling would oversimplify tumour boundaries. The adaptive approach selects the coarsest segmentation that still resolves the relevant tumour-discriminative structure.

\textbf{Step 3: Node Feature Extraction} - For each superpixel, compute the 15 features described previously by aggregating information from the 4 MRI modalities.

\textbf{Step 4: Edge Construction} - Build graph edges by determining spatial adjacency: two superpixel nodes are connected if they share a boundary in the image space. This creates a planar graph with average degree 4-6.

\textbf{Step 5: Node Label Assignment} - Assign binary labels (tumor/non-tumor) to each node based on majority voting: if $>$50\% of voxels within a superpixel are labeled as tumor in ground truth, the node is labeled positive.

\textbf{Computational Efficiency and Dimensionality Reduction:} Each 3D MRI volume spans $240 \times 240 \times 155 = 8{,}928{,}000$ voxels per modality. The superpixel-graph representation replaces this dense grid with a sparse set of nodes, one per brain-region superpixel per active slice. The theoretical compression ratio for the spatial representation can be expressed as:

\begin{equation}
R_{\text{theory}} = \frac{240 \times 240 \times 155 \times f_{\text{brain}}}{N_{\text{sp}} \times F}
\label{eq:compression_theory}
\end{equation}

where $f_{\text{brain}}$ is the fraction of non-zero (brain-masked) voxels, $N_{\text{sp}} = 200$ is the SLIC target superpixel count per slice (the \texttt{n\_segments} parameter), and $F = 15$ is the number of engineered features per node. With the often-cited assumption $f_{\text{brain}} \approx 0.30$:

\begin{equation}
R_{\text{theory}} = \frac{8{,}928{,}000 \times 0.30}{200 \times 15} = \frac{2{,}678{,}400}{3{,}000} \approx 893 \approx 890\times
\label{eq:compression_890}
\end{equation}

\textbf{Measured verification:} Empirical measurement across 50 BraTS 2021 patients (using the stored \texttt{brain\_mask} arrays in the preprocessed \texttt{.npz} files and saved per-patient graph \texttt{.pt} files) yields:

\begin{itemize}
\item Actual $f_{\text{brain}} = 0.154 \pm 0.017$ (not 0.30 — skull-stripped BraTS volumes are more compact than assumed)
\item Active slices per patient: $137.5 \pm 6.9$ out of 155 total
\item Actual superpixels per active slice: $46.4$ (SLIC adapts down from the 200 target for smaller brain regions)
\item Actual total nodes per patient: $3{,}909 \pm 821$
\end{itemize}

The \textit{spatial} compression — total voxels divided by actual node count — is therefore:

\begin{equation}
R_{\text{actual}} = \frac{8{,}928{,}000}{3{,}909} \approx 2{,}284\times
\label{eq:compression_actual}
\end{equation}

The commonly cited $\approx 890\times$ figure (Eq.~\ref{eq:compression_890}) is thus a \textbf{conservative lower bound}: it uses $f_{\text{brain}} = 0.30$ (twice the measured value) and $N_{\text{sp}} = 200$ (the unreduced SLIC target), while the real compression from empty-slice exclusion and adaptive superpixel reduction reaches $\sim$2,284$\times$. All claims in this thesis use the conservative 890$\times$ figure.

\subsubsection{GraphSAGE Architecture}

We employ GraphSAGE (Graph Sample and Aggregate), a message-passing neural network designed for inductive learning on graphs with node features.

\textbf{Core Mechanism:} GraphSAGE updates node representations by aggregating information from neighborhood nodes through learned transformations:

\begin{equation}
h_{\mathcal{N}(v)}^{(k)} = \frac{1}{|\mathcal{N}(v)|} \sum_{u \in \mathcal{N}(v)} h_u^{(k)}
\end{equation}

\begin{equation}
h_v^{(k+1)} = \sigma \left( W \cdot \left[ h_v^{(k)} \oplus h_{\mathcal{N}(v)}^{(k)} \right] \right)
\end{equation}

where $h_v^{(k)}$ is the representation of node $v$ at layer $k$, $\mathcal{N}(v)$ is the neighborhood of $v$, $h_{\mathcal{N}(v)}^{(k)}$ is the mean-aggregated neighborhood representation, $\oplus$ denotes concatenation, $W$ is a learnable weight matrix, and $\sigma$ is ReLU activation.

\textbf{Architecture Specification:}
\begin{itemize}
    \item \textbf{Input Layer:} 15 node features
    \item \textbf{Hidden Layers:} 5 GraphSAGE layers with 256 hidden dimensions each
    \item \textbf{Activation:} ReLU activation after each layer
    \item \textbf{Normalization:} Batch normalization for training stability
    \item \textbf{Output Layer:} Single neuron with sigmoid activation for binary classification
    \item \textbf{Total Parameters:} 439,041
\end{itemize}


Each layer performs message passing (aggregate neighbors → concatenate with self → linear transform → ReLU activation). The 5-layer depth provides 5-hop receptive field sufficient for tumor context aggregation, as validated through ablation studies showing no benefit from 6 layers.

\textbf{Design Rationale:} The 5-layer depth allows message propagation up to 5 hops in the graph, enabling each node to aggregate information from a wide spatial context. Ablation studies confirmed that 6 layers provide no additional benefit. The 256 hidden dimensions balance expressiveness with parameter efficiency.

\subsubsection{Training Protocol}

\textbf{Loss Function:} Binary Cross-Entropy with Logits loss (BCEWithLogitsLoss), which combines a sigmoid activation with binary cross-entropy in a numerically stable single operation:

\begin{equation}
\mathcal{L} = -\frac{1}{N} \sum_{i=1}^{N} \left[ y_i \log(\sigma(\hat{z}_i)) + (1-y_i) \log(1-\sigma(\hat{z}_i)) \right]
\end{equation}

where $\hat{z}_i$ is the raw logit output and $\sigma(\cdot)$ is the sigmoid function. Class imbalance (tumor nodes $\sim$10\%, non-tumor $\sim$90\%) is addressed through the Dice coefficient as the primary evaluation metric, which is insensitive to true-negative dominance.

\textbf{Optimization:} AdamW optimizer with learning rate 0.001, $\beta_1=0.9$, $\beta_2=0.999$, weight decay $0.01$.

\textbf{Learning Rate Schedule:} OneCycleLR scheduler with cosine annealing strategy and a warm-up phase comprising 10\% of total training steps, cycling from a low initial rate to the maximum learning rate of 0.001 and back.

\textbf{Batch Configuration:} Effective batch size 64, implemented as batch size 32 with 2 gradient accumulation steps. This configuration was chosen to fit within the RTX 2060 6\,GB VRAM constraint while maintaining stable gradient estimates. See Section~4.4 for ablation notes on batch size and the single-fold protocol used during architectural search.

\textbf{Training Duration:} 50 epochs per fold. The OneCycleLR scheduler provides structured learning rate progression across all training steps, with the cosine decay in the latter half providing implicit regularisation.

\textbf{Hardware:} NVIDIA RTX 2060 with 6GB VRAM, demonstrating accessibility on consumer-grade GPUs. Training time: $\sim$5 hours per fold.

\textbf{Reproducibility:} Deterministic training with manual seed 42 for Python, NumPy, PyTorch, and CUDA. Environment variables set: \texttt{CUBLAS\_WORKSPACE\_CONFIG=:4096:8}, \texttt{CUDNN\_DETERMINISTIC=True}.

\subsubsection{Ensemble Strategy}

To improve robustness and performance, we implement soft voting ensemble:

\begin{equation}
P_{ensemble}(tumor) = \frac{1}{5} \sum_{k=1}^{5} P_{fold\_k}(tumor)
\end{equation}

where $P_{fold\_k}(tumor)$ is the sigmoid output probability from fold $k$ model. Final prediction is tumor if $P_{ensemble} > 0.5$.

\textbf{Rationale:} Each fold trains on different 80\% patient subsets, learning complementary patterns. Averaging reduces prediction variance and improves generalization. Our experiments show consistent 2.5\% Dice improvement (90.39\% $\rightarrow$ 92.92\%).

\subsubsection{Evaluation Metrics}

\textbf{Primary Metric:} Dice Similarity Coefficient (DSC), the standard metric for medical image segmentation:

\begin{equation}
\text{Dice} = \frac{2 \times |A \cap B|}{|A| + |B|}
\end{equation}

where $A$ is predicted tumor region and $B$ is ground truth. Dice ranges from 0 (no overlap) to 1 (perfect agreement).

\textbf{Secondary Metrics:}
\begin{itemize}
    \item \textbf{Sensitivity (Recall):} $\frac{TP}{TP + FN}$ - ability to detect actual tumors
    \item \textbf{Specificity:} $\frac{TN}{TN + FP}$ - ability to correctly identify non-tumor regions
    \item \textbf{Precision:} $\frac{TP}{TP + FP}$ - accuracy of positive predictions
    \item \textbf{Inference Time:} Milliseconds per patient volume
    \item \textbf{GPU Memory:} Peak VRAM consumption during inference
    \item \textbf{Model Size:} Total trainable parameters
\end{itemize}


\subsection{Implementation Details}

The complete pipeline is implemented in Python using PyTorch and PyTorch Geometric frameworks. Medical image processing utilizes SimpleITK and nibabel for NIfTI file handling, while graph construction employs scikit-image's SLIC algorithm for superpixel generation. Preprocessed data is stored as compressed .npz files (4 modalities $\times$ 155 slices per BraTS volume) for efficient loading during training. Graph representations are structured as PyTorch Geometric Data objects containing node features (N $\times$ 15), edge indices (2 $\times$ E), and labels (N $\times$ 1).

Several implementation choices warrant emphasis for reproducibility: (1) \textbf{deterministic training} with manual seed 42 across Python, NumPy, PyTorch, and CUDA environments using \texttt{CUBLAS\_WORKSPACE\_CONFIG} and \texttt{CUDNN\_DETERMINISTIC} flags; (2) \textbf{effective batch configuration} using batch size 32 with 2 gradient accumulation steps (effective batch 64) to fit within 6\,GB VRAM while maintaining stable gradient estimates; (3) \textbf{comprehensive integrity auditing} with 15 automated validation checks detecting data leakage, patient overlap, and normalisation errors; (4) \textbf{efficient data structures} converting 8.9M-voxel volumes into $\sim$3,900-node graphs ($\geq$890$\times$ spatial compression; measured $\sim$2,284$\times$) while preserving tumour-discriminative features; and (5) \textbf{ensemble post-processing} using soft voting to aggregate sigmoid probabilities from all 5 folds prior to thresholding at 0.5. All experiments were conducted on consumer-grade hardware (NVIDIA RTX 2060, 6GB VRAM), demonstrating the accessibility of the proposed approach. The complete codebase with documentation is available for independent reproduction.

\vspace{5mm}

This chapter has presented a comprehensive and rigorous methodology for graph-based brain tumour segmentation. The proposed framework systematically transforms 3D MRI volumes into graph representations through superpixel-based node generation and feature engineering designed to avoid ground-truth leakage. The GraphSAGE architecture with 5 layers and 256 hidden dimensions achieves an optimal balance between expressiveness and parameter efficiency, as validated through systematic ablation studies. Patient-level stratified cross-validation ensures statistically valid performance estimation, while the ensemble strategy leverages model diversity for robust predictions. Comprehensive integrity validation with 15 auditing checks — covering data isolation, feature validity, graph structure, reproducibility, and output verification — ensures full reproducibility and eliminates data leakage, addressing critical concerns in medical AI research. The modular pipeline design facilitates efficient experimentation, adaptation to new datasets, and potential clinical translation. The subsequent chapter presents experimental results demonstrating that this methodology achieves competitive accuracy (92.92\% Dice) with substantial efficiency improvements (6.9$\times$ inference speedup, 156$\times$ parameter reduction) suitable for real-world clinical deployment.

\clearpage

